\paragraph{A precise limitation of the norm-ideal Massey formula.}
Sharifi's \emph{Massey products and ideal class groups},
J.~reine angew.~Math.~603 (2007), 1--33, has an author
version dated 26 July 2006 at
\url{https://www.math.ucla.edu/~sharifi/publ/massey.pdf}.
The page numbers in this paragraph refer to that forty-page
author version.  Section~4, pp.~17--18, fixes products with
consecutive copies of one Kummer class followed by a final
class.  The proper defining systems on pp.~21--22 require
a norm equation and an additional power condition.
Theorem~4.3 on p.~22 expresses the resulting value by
transgression, modulo the specified lower-product
indeterminacy.  Theorem~5.4 on p.~26, with its full proof
on pp.~26--27, gives the corresponding norm-ideal formula
for the prescribed ramification set.

These exact definitions, theorem statements and the complete
Theorem~5.4 proof were read for the present application.
Restricting the cocycles (CM.2) to $E(\mu_p)$ does trivialize
their coefficient characters.  It does not identify the
ordinary-normalized defining cochains (CM.19) with one of
those proper norm-based systems.  Nor has a reordering of
the repeated endpoint product $(u,v,u)$ into the consecutive
repeat product been proved to preserve its local primitives.
Even after such an identification, the norm-ideal formula
would require an arithmetic evaluation; it does not assert
that the ideal class is automatically zero or nonzero.
Thus this theorem is a possible evaluation method after
additional data are constructed, not a vanishing or
nonvanishing input to (CM.6).  No reading of the entire
forty-page paper or reconstruction of its foundational
proofs is claimed here.
